\section{Discussion and Outlook}
\label{sec:discussion}
In this section the results of the Metropolis algorithm and Wolff algorithm are compared and discussed. The results of the Metropolis algorithm are shown in \autoref{fig:Metropolis_Observables} and the results of the Wolff algorithm are shown in \autoref{fig:Wolff_Observables}. The single spin updates of the Metropolis algorithm are heavily affected by critical slowing down, resulting in a larger correlation time and thus higher variance of the measurements. This is further confirmed by \autoref{fig:Autocorrelation_Time_Metropolis}, where we see that the autocorrelation time $\tau$ increases significantly around $T_c$ for all lattice sizes, especially for larger lattice sizes. This is a clear indication of critical slowing down in the Metropolis algorithm, which is not present in the Wolff algorithm due to its cluster updates. Generally the more measurements are taken, the more accurate the results will be, however this also increases the computational time. The Wolff algorithm allows for more measurements to be taken in a given time frame due to its reduced correlation time.
Finer and wider temperature grid, especially around $T_c$ would enable us to better capture the behavior of the system around the critical temperature, and thus more accurately determine $T_c$ and the critical exponents. You could further improve the accuracy by estimating $\tau$ between every measurement, and only record measurements that are statistically independent, which would further reduce the variance of the measurements. However this would also increase the computational time, especially for the Metropolis algorithm, where $\tau$ is significantly larger around $T_c$. The Wolff algorithm would benefit less from this improvement, as its correlation time is already significantly reduced compared to the Metropolis algorithm. For further improvements a combination of CUDA and Pytorch could be used to further speed up the simulations, especially for larger lattice sizes. This would allow for more measurements to be taken in a given time frame, thus improving the accuracy of the results. Overall, the Wolff algorithm is more efficient and accurate for studying the XY model and its phase transition compared to the Metropolis algorithm, especially near the critical temperature where critical slowing down is most pronounced.